\newcommand{\CLASSINPUTinnersidemargin}{15mm}
\documentclass[journal]{IEEEtran}

\usepackage{graphicx}
\usepackage{siunitx}
\usepackage{bm}
\usepackage{cite}
\usepackage{flushend}
\usepackage{svg}

% Usage: \svg{Path}{Caption}{Width}
\newcommand{\svg}[3]{
    \begin{figure}[!h]
        \centering
        \includesvg[width=#3\linewidth]{#1}
        \caption{#2}
    \end{figure}
}

\begin{document}

\markboth{Analog VLSI Design NTUA Course}{}

\title{A Comparative Analysis of Capacitor-Less LDO Topologies in 130nm CMOS Technology}
\author{Kourosh N. Baygi, \IEEEmembership{Student Member}, \IEEEmembership{IEEE}, \textit{National Technical University of Athens}, \textit{Greece}, \textit{School of Electrical \& Computer Engineering}, el22174@mail.ntua.gr \vspace{-14pt}}
\maketitle

\begin{abstract}
This paper presents the design and comparison of three output-capacitor-less (OCL) PMOS low dropout (LDO) regulators implemented in a 130nm CMOS technology. Each LDO topology was designed using a different OTA (5T OTA, Current Mirror OTA, Telescopic OTA) while the PMOS Pass Device was kept identical. All three LDO configurations drive a 50 pF load, deliver a load current of {\bfseries \SI[per-mode=symbol,detect-weight]{500}{\micro\ampere}}, support programmable output voltages ({\bfseries \SI[per-mode=symbol,detect-weight]{2.3}{\volt}}, {\bfseries \SI[per-mode=symbol,detect-weight]{2.1}{\volt}} and {\bfseries \SI[per-mode=symbol,detect-weight]{1.6}{\volt}}), and have a consumption strictly below {\bfseries \SI[per-mode=symbol,detect-weight]{100}{\micro\ampere}}. Simulation results demonstrate that the 5T OTA provides the most balanced final results, whereas the Telescopic OTA maximizes loop gain and power supply rejection ratio (PSRR) at the sever expense of the Error Amplifier's output voltage swing. Furthermore, the Current Mirror OTA is shown to be incapable of satisfying the target specifications with the standardized PMOS pass device sizing.
\end{abstract}

\section{Introduction and Theoretical Framework}
\IEEEPARstart{L}{}ow dropout (LDO) voltage regulators are one of the most essential building blocks in power-management systems. Specially with today's increasing demand for system-on-chip solutions, many mixed-signal systems utilize LDO regulators to deliver local supply voltages with decreased noise.

This work will focus on desiging three LDOs. Each LDO topology must meet the proposed specifications that are presented in Table I. The general LDO schematic is also shown in Figure 1.

\begin{table}[!ht]
\renewcommand{\arraystretch}{1.4}
\caption{Technical LDO Target Specifications}
\centering
\begin{tabular}{|l|l|}
\hline
\textbf{Parameter} & \textbf{Proposed Value}\\
\hline
Power Supply Voltage & \qty{3}{\volt}\\
\hline
Bias Current $I_{bias}$ & \qty{5}{\micro\ampere}\\
\hline
Load Current $I_{L}$ & \qty{500}{\micro\ampere}\\
\hline
Load Capacitance $C_{L}$ & \qty{50}{\pico\farad}\\
\hline
Reference Voltage $V_{REF}$ & {$\qty{1.2}{\volt} \pm \qty{1}{\percent}$}\\
\hline
Output Voltage $V_{OUT}$ & \qty{2.3}{\volt}\\
\hline
Trim Voltages & \qty{2.1}{\volt}, \qty{1.6}{\volt}\\
\hline
Compensation Capacitor $C_c$ & \qty{<7}{\pico\farad}\\
\hline
Phase Margin @ No-Load & \qty{>60}{\degree}\\
\hline
Phase Margin @ Full-Load & \qty{>70}{\degree}\\
\hline
PSRR @ \qty{1}{\hertz} & \qty{>65}{\decibel}\\
\hline
PSRR @ \qty{1}{\mega\hertz} & \qty{>2}{\decibel}\\
\hline
PSRR @ \qty{10}{\mega\hertz} & \qty{>2}{\decibel}\\
\hline
Output Referred Noise @ \qty{10}{\hertz} & \qty{<-115}{\decibel\volt\slash\sqrt{\Hz}}\\
\hline
Output Referred Noise @ \qty{1}{\kilo\hertz} & \qty{<-125}{\decibel\volt\slash\sqrt{\Hz}}\\
\hline
Output Referred Noise @ \qty{10}{\kilo\hertz} & \qty{<-134}{\decibel\volt\slash\sqrt{\Hz}}\\
\hline
Output Referred Noise @ \qty{100}{\kilo\hertz} & \qty{<-143}{\decibel\volt\slash\sqrt{\Hz}}\\
\hline
Output Referred Noise @ \qty{1}{\mega\hertz} & \qty{<-155}{\decibel\volt\slash\sqrt{\Hz}}\\
\hline
\end{tabular}
\end{table}

\svg{images/pmos_ldo.svg}{General LDO Schematic}{1}

\subsection{Stability}

In this work we only use one-stage error amplifiers. A small signal model of the LDO schematic can be seen in Figure 2 where the error amplifier's transconductance, output resistance and output load are represented by $g_{m,EA}$, $R_1$ and $C_1$ respectively. Similarly, $g_{m,p}$, $R_2$ and $C_2$ represent the same values for the PMOS pass device. The LDO's transfer function will have two poles $\omega_{p1}$ (dominant) and $\omega_{p2}$ (non-dominant) and a zero $\omega_{z}$ all of which can be approximated with:

\begin{equation}
    \omega_{p1} = \frac{1}{g_{m,p} R_1 R_2 C_c}
\end{equation}

\begin{equation}
    \omega_{p2} = \frac{g_{m,p}}{C_1 + C_2}
\end{equation}

\begin{equation}
    \omega_{z} = \frac{-1}{C_c(1 \slash g_{m,p} - R_c)}
\end{equation}

\svg{images/pmos_ldo_ss.svg}{Small Signal Model of LDO}{1}

\nocite{*}
\bibliographystyle{IEEEtran}
\bibliography{IEEEabrv,report}

\begin{IEEEbiographynophoto}
{Kourosh N. Baygi} was born in Mashhad, Iran, in 2003 and is currently an undergraduate student in the School of Electrical \& Computer Engineering in National Technical University of Athens.
\end{IEEEbiographynophoto}

\end{document}